\chapter{Introduction}\label{chap:introduction}
Railway networks play an essential role in modern transportation systems. They support the daily mobility of passengers, 
contribute to freight transportation, and constitute a strategic infrastructure for sustainable mobility. In dense railway 
networks such as the Belgian one, the ability to operate trains efficiently, safely, and reliably is particularly 
important. A large number of trains, stations, tracks, and operational constraints interact continuously, making railway 
traffic management a complex task.

This complexity is not only due to the size of the infrastructure, but also to the nature of railway operations. Trains 
must follow predefined routes and timetables, share limited track capacity, respect safety constraints, and interact with 
infrastructure elements such as switches, platforms, and stations. A small disruption in one part of the network can 
propagate to other trains and stations, affecting punctuality and service quality. As a result, railway operators and 
researchers increasingly rely on digital tools to better understand, simulate, and analyse railway systems.

In this context, the concept of a digital twin has gained increasing attention. Informally, a digital twin can be 
understood as a virtual counterpart of a real physical system. It is not only a static model or a simple visualization of 
the system, but rather a digital representation that can be connected to data, updated over time, and used to observe, 
simulate, and analyse the behaviour of the physical system. Applied to railways, a digital twin may represent parts of 
the railway infrastructure, train movements, operational constraints, and traffic conditions. Such a representation can 
help explore how the network behaves, test different scenarios, and support decision-making.

Simulation is a central component of this vision. Before a railway digital twin can be used for prediction, optimization, 
or real-time decision support, it must first be possible to construct a coherent simulation model from available 
infrastructure and operational data. Simulation makes it possible to reproduce train movements in a controlled 
environment, evaluate whether generated routes are feasible, and compare simulated behaviour with real railway data.

This thesis focuses on the construction of a simulation-oriented railway digital twin using SUMO.\@ SUMO is a traffic 
simulation platform originally designed for road traffic, but its flexible network and route representation make it 
possible to adapt it to railway use cases. The work investigates how railway infrastructure and operational data can be 
transformed into simulation-ready artefacts, and how different levels of railway network representation can be used to 
support train simulation.


\section{Problem Statement}
Despite the availability of open data sources and simulation platforms, constructing a realistic railway simulation 
remains a difficult task. Railway data is heterogeneous by nature: infrastructure data describes physical components 
such as stations, tracks, switches, and platforms, while operational data describes train movements, schedules, and 
punctuality records. These datasets do not always share the same structure, granularity, or level of detail. Consequently, 
integrating them into a single simulation framework requires several processing, modeling, and approximation steps.

A first difficulty concerns the representation of the railway infrastructure. At a high level, a railway network can be 
represented as a graph where stations are connected by railway links. This abstraction is useful because it simplifies 
the network and makes it easier to reconstruct train movements between stations. However, it does not represent detailed 
infrastructure elements such as platforms, switches, parallel tracks, or station-level routing constraints. These details 
are important when the objective is to simulate train circulation more realistically.

A second difficulty appears when moving toward a more detailed representation of the network. At the micro level, tracks, 
platforms, and switches must be represented explicitly. This provides a more realistic description of the infrastructure, 
but it also introduces new challenges. The network must be topologically consistent, platforms must be associated with 
the correct tracks, and valid routes must exist between consecutive stations. Even small errors in the reconstructed 
infrastructure can prevent trains from being routed correctly.

A third difficulty concerns the integration of real operational data. Punctuality data provides information about train 
stops, arrival times, departure times, and delays, but it does not directly provide the exact physical route followed by 
each train inside the infrastructure. Therefore, a route must be reconstructed from the available station sequences and 
mapped onto the generated simulation network. When the network is simplified, this process is easier but less precise. 
When the network is detailed, the process becomes more realistic but also more fragile.

These challenges lead to the central research question of this thesis:

\begin{quote}
    How can a railway digital twin simulation be constructed using SUMO and heterogeneous railway datasets, and to what 
    extent can real operational data be integrated into feasible macro-level and micro-level railway simulations?
\end{quote}

This question is addressed by studying two complementary modeling approaches. The first one is a macro-level model, where 
stations are represented as nodes and station-to-station connections as edges. The second one is a micro-level model, 
where detailed infrastructure elements such as tracks, switches, and platforms are explicitly represented. Together, these 
two approaches make it possible to evaluate the trade-off between simplicity, scalability, infrastructure realism, and 
simulation feasibility.

\section{Objectives}
The main objective of this thesis is to design and implement a data-driven pipeline for constructing a simulation-oriented 
digital twin of a railway network. The focus is not to build a fully operational real-time digital twin, but rather to 
study the foundations required to generate, simulate, and evaluate railway traffic scenarios from real-world data.

More specifically, this work aims to construct railway simulation models at two different levels of abstraction. The 
macro-level model is designed to represent the railway network as a station-to-station graph. Its purpose is to provide a 
simplified but functional simulation environment in which real train trajectories can be reconstructed from punctuality 
data. This model serves as a first validation of the complete pipeline, from data preprocessing to simulation execution.

The second objective is to develop a micro-level model that represents the infrastructure in greater detail. This includes 
the processing of track data, the detection of switches, the identification of platforms, the segmentation of track 
sections, and the construction of routing structures. The objective of this model is to move closer to the physical 
structure of the railway network and to evaluate whether more detailed simulations can be produced from the available data.

Another objective is to compare different trip generation strategies. In the macro-level model, the focus is mainly on the 
use of real operational data to reconstruct train movements. In the micro-level model, two approaches are considered: 
synthetic trip generation and real-data-based trip generation. Synthetic trips are used to test the feasibility and 
connectivity of the generated infrastructure, while real-data trips aim to reproduce observed railway operations more 
accurately.

Finally, this thesis aims to identify the main limitations encountered when building a railway digital twin simulation 
from heterogeneous data. These limitations concern data availability, infrastructure completeness, modeling 
assumptions, routing feasibility, and constraints imposed by the simulation platform. By analysing these issues, the work 
provides methodological insights for future developments of railway digital twins.

The expected outcome of this thesis is therefore not a perfect reproduction of the Belgian railway network, but a 
structured and evaluated prototype demonstrating how railway data can be transformed into SUMO-compatible simulation 
models. The contribution lies in the development of the modeling pipeline, the comparison between macro-level and 
micro-level representations, and the identification of the technical challenges that must be addressed to progress toward 
more complete railway digital twins.

\section{Thesis Structure}
The remainder of this thesis is organized as follows.

Chapter~\ref{chap:stateoftheart} presents the state of the art. It introduces the concept of digital twins, first from a 
general perspective and then in the context of railway systems. It also reviews the main components of railway 
infrastructure and operations, discusses common approaches for railway network representation, and presents existing 
railway simulation tools. Particular attention is given to SUMO and to its suitability for railway simulation.

Chapter~\ref{chap:methodology} describes the proposed methodology. It positions the work with respect to the research 
gaps identified in the literature and presents the overall framework followed throughout the project. This chapter also 
justifies the choice of SUMO as the simulation platform and describes the infrastructure and operational datasets used in 
the implementation.

Chapter~\ref{chap:implementation} details the implementation of the proposed approach. It is divided into two main parts. 
The first part presents the macro-level model, including data processing, network generation, trip generation, simulation, 
and its limitations. The second part presents the micro-level model, including infrastructure modelling, switch detection, 
platform detection, track segmentation, dual graph construction, trip generation, simulation, and its limitations.

Chapter~\ref{chap:results} presents the results obtained from the implemented models. It analyses dataset coverage, 
infrastructure modelling quality, simulation results, and computational performance. The objective of this chapter is to 
evaluate how well the generated models represent the available railway data and how effectively they can be used for 
simulation.

Chapter~\ref{chap:limitations} discusses the main limitations of the proposed methodology. While some limitations are 
briefly introduced in the implementation chapter, this chapter provides a more detailed analysis of issues related to 
data quality, modelling assumptions, SUMO constraints, and the generalization of the approach to other railway networks.

Chapter~\ref{chap:conclusion} concludes the thesis by summarizing the main contributions and findings. It also discusses 
the practical implications of the work and presents future research directions, including counterfactual simulation, track 
work planning, real-time forecasting, incident propagation, and the generalization of the proposed framework.